% ═══════════════════════════════════════════════════════════════
% ML-07c: Musterlösung - Querer y Poder
% Abgeleitet aus LE-07c (Abschnitt 9: Lösungen)
% ═══════════════════════════════════════════════════════════════

\documentclass[11pt, a4paper]{article}

\usepackage[utf8]{inputenc}
\usepackage[T1]{fontenc}
\usepackage[ngerman]{babel}
\usepackage{helvet}
\renewcommand{\familydefault}{\sfdefault}

\usepackage[margin=2cm]{geometry}
\usepackage{enumitem}
\usepackage{tabularx}
\usepackage{booktabs}

\usepackage{xcolor}
\usepackage{tcolorbox}

\usepackage{fancyhdr}
\usepackage{lastpage}

% Farben
\definecolor{spanisch}{HTML}{c41e3a}
\definecolor{grau}{HTML}{888888}
\definecolor{hellgrau}{HTML}{f5f5f5}
\definecolor{orange}{HTML}{b35c00}
\definecolor{loesung}{HTML}{c62828}

\newcommand{\fachfarbe}{spanisch}

% Variablen
\newcommand{\thema}{Querer y Poder -- Modalverben beim Einkaufen}
\newcommand{\sequenz}{07}
\newcommand{\block}{c}
\newcommand{\fach}{Spanisch}
\newcommand{\klasse}{7}
\newcommand{\maxpunkte}{24}

% Kopf-/Fußzeile
\pagestyle{fancy}
\fancyhf{}
\renewcommand{\headrulewidth}{0.4pt}
\renewcommand{\footrulewidth}{0.4pt}
\lhead{\fach{} -- Klasse \klasse}
\chead{\textcolor{loesung}{\textbf{ML-\sequenz\block{} (MUSTERLÖSUNG)}}}
\rhead{}
\lfoot{}
\rfoot{Seite \thepage\ von \pageref{LastPage}}

% Befehle
\newcommand{\punktebox}[1]{\hfill\fbox{\small \_/#1}}
\newcommand{\lsg}[1]{\textcolor{loesung}{\textbf{#1}}}
\newcommand{\luecke}[1]{\rule{#1}{0.4pt}}
\newcommand{\es}[1]{\textcolor{\fachfarbe}{\textbf{#1}}}

% Merkkasten
\newtcolorbox{merkkasten}[1][]{
    colback=hellgrau,
    colframe=\fachfarbe,
    boxrule=1pt,
    arc=0pt,
    left=8pt, right=8pt, top=8pt, bottom=8pt,
    fonttitle=\bfseries\color{\fachfarbe},
    title=#1
}

% Hinweis
\newtcolorbox{hinweis}{
    colback=white,
    colframe=orange,
    boxrule=1pt,
    arc=0pt,
    left=5pt, right=5pt, top=5pt, bottom=5pt,
    fonttitle=\bfseries,
    title=Hinweis
}

% Aufgabe
\newenvironment{aufgabe}[1]{%
    \vspace{10pt}
    \noindent\textbf{\textcolor{\fachfarbe}{Aufgabe #1}}
    \vspace{5pt}
    \begin{list}{}{\leftmargin=0pt}
    \item[]
}{%
    \end{list}
}

% ═══════════════════════════════════════════════════════════════
\begin{document}

% Kopfbereich
\begin{center}
    {\Large\bfseries\textcolor{\fachfarbe}{Arbeitsblatt: \thema}}\\[0.3em]
    {\large\textcolor{loesung}{\textbf{--- MUSTERLÖSUNG ---}}}
\end{center}

\vspace{5pt}

\noindent
\begin{tabularx}{\textwidth}{|X|X|X|}
\hline
\textbf{Name:} \textcolor{loesung}{Musterschüler/in} &
\textbf{Klasse:} \klasse &
\textbf{Datum:} \textcolor{loesung}{XX.XX.XXXX} \\
\hline
\end{tabularx}

\vspace{15pt}

% ═══════════════════════════════════════════════════════════════
% MERKKASTEN 1 (mit Lösungen)
% ═══════════════════════════════════════════════════════════════

\begin{merkkasten}[Merkkasten: Das Verb \textit{querer} (wollen) -- Stammwechsel e $\rightarrow$ ie]

\begin{center}
\begin{tabular}{|l|l|l|l|}
\hline
\textbf{Person} & \textbf{Pronomen} & \textbf{querer} & \textbf{Stammwechsel} \\
\hline
1. Sg. & yo & \lsg{quiero} & e $\rightarrow$ ie \\
\hline
2. Sg. & tú & \lsg{quieres} & e $\rightarrow$ ie \\
\hline
3. Sg. & él/ella & \lsg{quiere} & e $\rightarrow$ ie \\
\hline
1. Pl. & nosotros/-as & \lsg{queremos} & -- \\
\hline
\end{tabular}
\end{center}

\end{merkkasten}

\vspace{10pt}

% ═══════════════════════════════════════════════════════════════
% MERKKASTEN 2 (mit Lösungen)
% ═══════════════════════════════════════════════════════════════

\begin{merkkasten}[Merkkasten: Das Verb \textit{poder} (können) -- Stammwechsel o $\rightarrow$ ue]

\begin{center}
\begin{tabular}{|l|l|l|l|}
\hline
\textbf{Person} & \textbf{Pronomen} & \textbf{poder} & \textbf{Stammwechsel} \\
\hline
1. Sg. & yo & \lsg{puedo} & o $\rightarrow$ ue \\
\hline
2. Sg. & tú & \lsg{puedes} & o $\rightarrow$ ue \\
\hline
3. Sg. & él/ella & \lsg{puede} & o $\rightarrow$ ue \\
\hline
1. Pl. & nosotros/-as & \lsg{podemos} & -- \\
\hline
\end{tabular}
\end{center}

\end{merkkasten}

\vspace{15pt}

% ═══════════════════════════════════════════════════════════════
% AUFGABE 1 (mit Lösungen)
% ═══════════════════════════════════════════════════════════════

\begin{aufgabe}{1}
Verbinde die Pronomen mit der richtigen Form. \punktebox{6}
\end{aufgabe}

\begin{center}
\begin{tabular}{rcl}
\textbf{querer:} & & \\[0.5em]
yo & \lsg{------} & \lsg{quiero} \\[0.8em]
tú & \lsg{------} & \lsg{quieres} \\[0.8em]
él/ella & \lsg{------} & \lsg{quiere} \\[1em]
\textbf{poder:} & & \\[0.5em]
yo & \lsg{------} & \lsg{puedo} \\[0.8em]
tú & \lsg{------} & \lsg{puedes} \\[0.8em]
él/ella & \lsg{------} & \lsg{puede} \\
\end{tabular}
\end{center}

\textit{\textcolor{grau}{Je 1P pro richtige Zuordnung}}

\vspace{15pt}

% ═══════════════════════════════════════════════════════════════
% AUFGABE 2 (mit Lösungen)
% ═══════════════════════════════════════════════════════════════

\begin{aufgabe}{2}
Setze die richtige Form von \textit{querer} oder \textit{poder} ein. \punktebox{6}
\end{aufgabe}

\begin{enumerate}[label=\alph*)]
    \item Yo \lsg{quiero} comprar una camiseta roja. (wollen)
    \item ¿\lsg{Puedo} ver esa falda azul? (können, ich)
    \item Él \lsg{quiere} unos pantalones negros. (wollen)
    \item Nosotros \lsg{podemos} pagar con tarjeta. (können)
    \item ¿Tú \lsg{quieres} probarte los zapatos? (wollen)
    \item María no \lsg{puede} encontrar su talla. (können)
\end{enumerate}

\textit{\textcolor{grau}{Je 1P pro richtige Form}}

\vspace{15pt}

% ═══════════════════════════════════════════════════════════════
% AUFGABE 3 (mit Lösungen)
% ═══════════════════════════════════════════════════════════════

\begin{aufgabe}{3}
Vervollständige den Einkaufsdialog. \punktebox{6}
\end{aufgabe}

\noindent
\textbf{Vendedor:} \es{¡Buenos días! ¿En qué puedo ayudarle?}

\noindent
\textbf{Cliente:} \lsg{Quiero una camiseta.} (ich will ein T-Shirt)

\noindent
\textbf{Vendedor:} \es{¿De qué color?}

\noindent
\textbf{Cliente:} \es{Rojo, por favor.} \lsg{¿Puedo probármela?} (kann ich es anprobieren?)

\noindent
\textbf{Vendedor:} \es{¡Claro! El probador está allí.}

\noindent
\textbf{Cliente:} \es{Me queda bien.} \lsg{¿Cuánto cuesta?} (wie viel kostet es?)

\noindent
\textbf{Vendedor:} \es{Son 25 euros.}

\noindent
\textbf{Cliente:} \lsg{¿Puedo pagar con tarjeta?} (kann ich mit Karte bezahlen?)

\noindent
\textbf{Vendedor:} \es{Sí, por supuesto.}

\vspace{0.5em}
\textit{\textcolor{grau}{Je 1,5P pro sinnvolle Ergänzung, max. 6P}}

\vspace{15pt}

% ═══════════════════════════════════════════════════════════════
% AUFGABE 4 (Beispiellösung)
% ═══════════════════════════════════════════════════════════════

\begin{aufgabe}{4}
Schreibe einen eigenen Einkaufsdialog (mind. 8 Zeilen). \punktebox{6}
\end{aufgabe}

\textbf{Beispiellösung:}

\lsg{Vendedor: ¡Hola! ¿Qué quieres?}

\lsg{Cliente: Quiero unos pantalones.}

\lsg{Vendedor: ¿De qué color?}

\lsg{Cliente: Azul. ¿Puedo ver esos?}

\lsg{Vendedor: Sí, claro. ¿Qué talla quieres?}

\lsg{Cliente: Talla 38. ¿Puedo probármelos?}

\lsg{Vendedor: Sí, el probador está allí. ¿Cuánto cuestan?}

\lsg{Vendedor: Son 40 euros.}

\lsg{Cliente: ¿Puedo pagar con tarjeta?}

\lsg{Vendedor: Sí, por supuesto. ¡Hasta luego!}

\vspace{0.5em}
\textit{\textcolor{grau}{Bewertung: Begrüßung (1P), querer verwenden (1P), poder verwenden (1P), Preis/Bezahlung (1P), Länge 8+ Zeilen (1P), Korrektheit (1P)}}

\vspace{10pt}

% ═══════════════════════════════════════════════════════════════
% BONUSBOX (mit Lösungen)
% ═══════════════════════════════════════════════════════════════

\begin{merkkasten}[Bonus: Die Stammwechsel-Regel]
\textbf{Wann wechselt der Stammvokal?}

\lsg{Bei betonter Silbe wechselt der Stammvokal. Bei nosotros/vosotros ist die Endung betont, daher kein Wechsel.}

\vspace{0.5em}
\textbf{Weitere Verben mit Stammwechsel:}

e $\rightarrow$ ie: \lsg{pensar (pienso), entender (entiendo), preferir (prefiero)}

o $\rightarrow$ ue: \lsg{dormir (duermo), volver (vuelvo), encontrar (encuentro)}
\end{merkkasten}

\vspace{20pt}
\hrule
\vspace{10pt}

\begin{flushright}
\textbf{Punkte:} \lsg{\maxpunkte} / \maxpunkte
\end{flushright}

\end{document}
